\documentclass[a4paper]{article}
% Español (México) con XeLaTeX: fontspec para fuentes Unicode, polyglossia
% para la división silábica y los nombres del documento en la variante mexicana.
\usepackage{fontspec}
\usepackage{polyglossia}
\setdefaultlanguage[variant=mexican]{spanish}
% Los nombres chinos van como «pinyin (original)»: xeCJK da a los caracteres
% Han su propia fuente; sin ella XeLaTeX los omite con un aviso.
\usepackage{xeCJK}
\setCJKmainfont{FandolSong-Regular.otf}
% polyglossia no traduce los nombres de listings.
\AtBeginDocument{\providecommand{\lstlistingname}{}\renewcommand{\lstlistingname}{Listado}%
  \providecommand{\lstlistlistingname}{}\renewcommand{\lstlistlistingname}{Índice de listados}}
\usepackage{amsmath, amssymb}
\usepackage{graphicx}
\usepackage[margin=2.5cm]{geometry}
\usepackage[most]{tcolorbox}
\usepackage{etoolbox}
\usepackage{listings}
\usepackage{hyperref}
\usepackage{booktabs}
\usepackage{subcaption}
\usepackage{float}

\newtcolorbox{knowledgebox}[1]{enhanced, colback=blue!5!white, colframe=blue!75!black, colbacktitle=blue!75!black, coltitle=white, fonttitle=\bfseries, title={#1}, attach boxed title to top left={yshift=-2mm, xshift=2mm}, boxrule=1pt, sharp corners}
\newtcolorbox{importantbox}[1]{enhanced, colback=yellow!10!white, colframe=yellow!80!black, colbacktitle=yellow!80!black, coltitle=black, fonttitle=\bfseries, title={#1}, sharp corners}
\newtcolorbox{warningbox}[1]{enhanced, colback=red!5!white, colframe=red!75!black, colbacktitle=red!75!black, coltitle=white, fonttitle=\bfseries, title={#1}, sharp corners}

\lstset{language=Python, basicstyle=\ttfamily\small, keywordstyle=\color{blue}, stringstyle=\color{red!60!black}, commentstyle=\color{green!60!black}, breaklines=true, frame=single, numbers=left, numberstyle=\tiny\color{gray}, captionpos=b, extendedchars=true}

\newcommand{\notetitle}{LlamaIndex RAG: principios y análisis del código fuente}
\newcommand{\noteauthors}{Elaboradas a partir de materiales públicos del curso}
\newcommand{\notedate}{2025}
\newcommand{\videochannel}{Wudaokou Nash (五道口纳什)}
\newcommand{\videopublishdate}{2025}
\newcommand{\videoduration}{aproximadamente 22 minutos}
\newcommand{\videourl}{https://www.bilibili.com/video/BV1Vm42157PL}
\newcommand{\videocoverpath}{cover.jpg}
\newcommand{\repourl}{https://github.com/hqhq1025/ai-course-notes}


% Footer with repo info
\usepackage{fancyhdr}
\pagestyle{fancy}
\fancyhf{}
\fancyfoot[C]{\thepage}
\fancyfoot[R]{\footnotesize\color{gray}\href{\repourl}{AI Course Notes}}
\renewcommand{\headrulewidth}{0pt}
\renewcommand{\footrulewidth}{0pt}

\begin{document}
\begin{titlepage}
\centering
{\Large Notas del curso\par}\vspace{1.2cm}
{\huge\bfseries \notetitle\par}\vspace{0.8cm}
{\large \noteauthors\par}\vspace{0.3cm}
{\large \notedate\par}\vspace{1.2cm}
\ifdefempty{\videocoverpath}{}{\includegraphics[width=0.82\textwidth,height=0.45\textheight,keepaspectratio]{\videocoverpath}\par}
\vfill
\begin{tcolorbox}[width=0.9\textwidth, colback=black!2!white, colframe=black!60, sharp corners]
\textbf{Autor o canal del video}: \videochannel\par\textbf{Fecha de publicación}: \videopublishdate\par\textbf{Duración del video}: \videoduration\par
\ifdefempty{\videourl}{}{\textbf{Enlace del video}: \href{\videourl}{\nolinkurl{\videourl}}\par}\par
\textbf{Página del proyecto}: \href{\repourl}{\nolinkurl{\repourl}}\par
\end{tcolorbox}
\end{titlepage}
\tableofcontents\newpage

\section{Conceptos básicos de RAG}

\subsection{Marco general de RAG}

RAG (Retrieval-Augmented Generation) es un marco moderno muy práctico y muy básico, indispensable tanto en la ingeniería como en la investigación.

\begin{importantbox}{Flujo central de RAG}
\begin{enumerate}
    \item El usuario emite una \textbf{Query} (pregunta)
    \item La Query se compara con el \textbf{Index} de los datos externos: los datos externos son datos que existen fuera del modelo de lenguaje, almacenados en disco (bases de datos estructuradas, documentos no estructurados, etc.)
    \item Los datos externos se someten a \textbf{Chunking} (división en fragmentos), y para cada Chunk se crea un \textbf{Embedding Vector}
    \item También se crea un Embedding Vector para la Query
    \item Se calcula la \textbf{similitud} (producto punto/coseno) entre la Query y cada Chunk, y se toman los Top-K Chunk más relevantes
    \item Los Chunk relevantes se usan como Context y se envían junto con la Query al LLM, que genera la respuesta
\end{enumerate}
\end{importantbox}

\subsection{Los dos Model centrales}

Un sistema RAG requiere dos modelos:
\begin{enumerate}
    \item \textbf{Embedding Model}: convierte el texto en una representación vectorial (por defecto text-embedding-ada-002, 1536 dimensiones)
    \item \textbf{Language Model}: genera la respuesta a partir del Context (por defecto GPT-3.5 Turbo)
\end{enumerate}

\begin{knowledgebox}{El valor central de RAG}
RAG exige que el LLM responda basándose únicamente en la Context Information recuperada, y no en el conocimiento que el propio modelo ya posee, lo que reduce las alucinaciones y aumenta la confiabilidad en escenarios con exigencias factuales altas.
\end{knowledgebox}

\subsection{Resumen de la sección}

La esencia de RAG es la «generación aumentada por recuperación»: primero se recupera de los datos externos la información más relevante, y después se hace que el LLM responda la pregunta basándose en esa información, lo que reduce las alucinaciones y aumenta la precisión.

\section{Análisis del código fuente de LlamaIndex: proceso de Indexing}

\subsection{Carga de documentos y Chunking}

\begin{enumerate}
    \item Se cargan todos los documentos del directorio Data (se obtiene una lista de Document)
    \item A cada documento se le aplica \textbf{Split}: el Chunk Size es de 1024 (en realidad alrededor de 997, restando los tokens de la ruta del documento), con un Overlap de 200
    \item El proceso de Split primero divide según sus propias reglas (produce alrededor de 759 splits) y luego los combina mediante \textbf{Merge} en los Chunk finales (alrededor de 22)
\end{enumerate}

\subsection{Creación de Node y Embedding}

\begin{enumerate}
    \item Cada Chunk crea un \textbf{Node} (con un node\_id único y el contenido de texto correspondiente)
    \item Se usa el Embedding Model (text-embedding-ada-002, con Batch Size 2048) para calcular el Embedding Vector de cada Node
    \item Al final se obtienen 22 Node, cada uno con un Embedding Vector de 1536 dimensiones ($512 \times 3$)
\end{enumerate}

\subsection{Resumen de la sección}

Proceso de Indexing: documento $\rightarrow$ Split $\rightarrow$ Merge $\rightarrow$ Chunk $\rightarrow$ Node $\rightarrow$ Embedding Vector. Cada paso tiene una implementación clara en el código fuente.
\section{Análisis del código fuente de LlamaIndex: flujo de Query}

\subsection{Proceso de recuperación}

\begin{enumerate}
    \item Se calcula el Embedding Vector del Query del usuario (también de 1536 dimensiones)
    \item Se calcula la similitud entre el Query Embedding y todos los Node Embedding
    \item Por defecto se toman los \textbf{Top-2} Node más relevantes
\end{enumerate}

El cálculo de similitud admite varias métricas: distancia euclidiana, producto punto, similitud coseno, entre otras.

\subsection{Prompt Template}

\begin{importantbox}{QA Template (extraído del código fuente)}
Plantilla central:
\begin{quote}
\textit{Context information is below. \{context\} Given the context information and not prior knowledge, answer the query. Query: \{query\}}
\end{quote}
Restricción clave: \textbf{responder únicamente con base en la Context Information, sin usar el conocimiento previo del modelo.}
\end{importantbox}

Además existen la Refine Template (que perfecciona una respuesta ya existente a partir de nuevo Context) y la Summarize Template, entre otras.

\subsection{Construcción del mensaje y llamada a la API}

El mensaje final construido incluye:
\begin{enumerate}
    \item \textbf{System Message}: definición de rol — "eres un experto en QA confiable, que siempre responde con base en la Context Information, sin usar conocimiento previo"
    \item \textbf{User Message}: incluye la Context Information (los documentos relevantes Top-K) + el Query
\end{enumerate}

El Response Mode por defecto es Compact (compacto).

\subsection{Resumen de la sección}

Flujo del Query: Query Embedding $\rightarrow$ coincidencia de similitud $\rightarrow$ Top-K Node $\rightarrow$ ensamblaje de la Prompt Template $\rightarrow$ el LLM genera la respuesta. Todo el proceso ocupa solo unas cinco líneas de código del framework, pero el código fuente interno es muy minucioso.

\section{Síntesis y ampliación}

\begin{enumerate}
    \item \textbf{RAG es una técnica muy básica y práctica}, LlamaIndex ofrece una implementación con un alto grado de encapsulamiento; a nivel de framework basta con unas cinco líneas de código para completar todo el proceso.
    \item \textbf{Entender el código fuente es la clave para dominar el framework}: LlamaIndex está encapsulado en un nivel muy alto, pero solo se domina de verdad al entender el proceso interno de Split/Merge/Node/Embedding/Retrieve/Template.
    \item \textbf{Componentes centrales}: el Embedding Model se encarga de la vectorización, el Language Model se encarga de la generación, el Prompt Template se encarga de organizar la entrada y el cálculo de similitud se encarga de la recuperación.
    \item \textbf{Direcciones de ampliación}: personalizar el Embedding Model, ajustar el Chunk Size y el Overlap, modificar el Response Mode, introducir estrategias de Retrieval más complejas, entre otras.
\end{enumerate}

\end{document}
